0.563 & 1.667\documentclass[11pt,a4paper]{article}
\usepackage[utf8]{inputenc}
\usepackage[T1]{fontenc}
\usepackage[margin=2.2cm]{geometry}
\usepackage{amsmath,amssymb,booktabs,tabularx,graphicx,xcolor,enumitem,hyperref}
\graphicspath{{figs/}}

\newcommand{\hint}[1]{\textcolor{gray}{\textit{[#1]}}}
\newcommand{\zref}{z_{\mathrm{ref}}}
\newcommand{\zLP}{z_{\mathrm{LP}}}
\setlength{\parindent}{0pt}\setlength{\parskip}{4pt}
\begin{document}

\begin{center}
{\LARGE\bfseries Final report}\\[4pt]
{\large Introduction to Optimization}
\end{center}

\begin{tabularx}{\textwidth}{@{}lX@{}}
\toprule
Team & T12 \\
Storyline, variant & B, 2 (cloud resource allocation, zones) \\
Members & Zhussupbekova Zere, Kabdulla Danial, Yermukhamed Tilesbay, Aitzhanov Tamerlan, Noken Rakhima \\
Repository & \url{https://github.com/zere90/12-checkpoint2} \\
Tag and commit hash & \texttt{final}, \hint{fill in after tagging} \\
\bottomrule
\end{tabularx}

\section{From the Checkpoint 1 model to the reference model}

Checkpoint 1 (22 September 2026) was our own, team-chosen formulation of storyline B, variant 2, on the same data (seed 12, variant 2) used ever since. The reference model built from Checkpoint 3 onward (\texttt{src/model/build.py:build\_b2\_model}, the model behind $\zref$, $\zLP$, Table 2 and Table 3 above) is not that formulation -- it is the model the course handout fixed for everyone from Checkpoint 3 on, so that every team's exact solver, LP relaxation and metaheuristics are comparable. This section compares the two models element by element, then checks what our \emph{own} Checkpoint 1 objective says about the solutions the reference model calls optimal.

\subsection*{Table 1. Element by element}
{\small
\begin{tabularx}{\textwidth}{@{}p{2.6cm}Xp{1.7cm}X@{}}
\toprule
Element & In Checkpoint~1 / in the reference model & Kept, changed, dropped & What it does to the meaning of the answer\\
\midrule
Variables & CP1: $x_{ij}\in\{0,1\}$, VM $i$ on server $j$ (350 vars, no auxiliaries). Reference: $y_{kl}\in\{0,1\}$ (same meaning, renamed) plus $t_l\ge0$, one per server (360 vars: 325 free $y$, 60 fixed to 0 by zones, 10 $t$). & Kept + added & $x_{ij}$'s meaning is unchanged; $t_l$ is new bookkeeping needed to turn a nonlinear per-server penalty into linear constraints (below), not a new decision about placement.\\
Objective & CP1: $\min \sum_j n_j(x)\,u_j(x)^2$ -- quadratic in utilization, weighted by how many VMs server $j$ hosts. Reference: $\min \sum_l t_l$ -- linear in the new variables, with $t_l$ forced by (B4) to sit above 4 tangent lines of $g(u)=u/(1-u)$. & Changed & Two changes at once: (i) the latency-vs-utilization curve itself changes (quadratic $\to$ piecewise-linear bound on an M/M/1 curve, see Assumption~1), and (ii) the reference model drops the per-server VM-count weighting $n_j(x)$ entirely -- it sums one latency number per server, not one per VM hosted. A solution that is best for ``total latency summed once per server'' need not be best for ``total latency experienced across all VMs,'' which is exactly what Section~3 below finds.\\
Constraint group 1 (complete assignment) & CP1 (C1): $\sum_j x_{ij}=1\ \forall i$. Reference (B1): $\sum_l y_{kl}=1\ \forall k$. & Kept & Identical meaning under the variable renaming; every VM still placed exactly once.\\
Constraint group 2 (capacity) & CP1 (C2)/(C3): $\sum_i x_{ij}\,\mathrm{cpu\_req}_i\le \mathrm{cpu\_cap}_j$, same for memory. Reference (B2)/(B3): same two inequalities in $y$. & Kept & Identical meaning; capacity and the zone restriction (C4/forbidden-pair upper bounds) and integrality (C5) all carry over unchanged in substance, only renamed.\\
Assumption 1 (latency--utilization function) & CP1 Section 5: we chose $f(u)=u^2$ \emph{specifically rejecting} the classical M/M/1 curve $g(u)=u/(1-u)$, because $g$ blows up as $u\to1$ and would make the objective unstable (Tilesbay's and Tamerlan's signed paragraphs). Reference model: uses $g(u)=u/(1-u)$ itself -- the very curve we rejected -- but replaced by a 4-segment piecewise-linear \emph{lower envelope} (tangents at $u=0,0.4,0.7,0.9$), which stays finite by construction. & Changed & The reference model keeps the accelerating-congestion shape our Tamerlan argued CPU contention should have, while avoiding the blow-up we were worried about, by bounding rather than evaluating $g$ exactly. So neither model is simply ``more correct'' than the other: CP1 chose boundedness over shape-fidelity; the reference model gets both, but only as a lower bound, never the true $g(u)$ value (CP3 S1, Tilesbay).\\
\bottomrule
\end{tabularx}}

\subsection*{Own objective on the reference solutions}
Both rows use instance \texttt{cp34\_instance(12,'B',2)} (seed 12, variant 2 -- the same data as Checkpoint 1). The reference-optimum assignment was recomputed by solving \texttt{build\_b2\_model(d)} exactly (HiGHS, proven, matches $\zref=5.268322$ to 6 decimals); the approximate-method assignment is SA with seed 0 (Table 2: all 10 SA seeds tie $\zref$ exactly, std $=0$, so any seed represents ``the'' SA solution). The two assignments are \emph{not} identical -- the reference model has multiple optimal vertices, as Checkpoint 3's Rakhima already noted for the LP relaxation (S2) -- which is itself part of the finding below.

\begin{tabularx}{\textwidth}{@{}Xrr@{}}
\toprule
Solution & Value under the reference objective & Value under my Checkpoint~1 objective\\
\midrule
Optimum of the reference model ($\zref$) & 5.268322 & 7.530112\\
Best solution of the approximate methods (SA, seed 0) & 5.268322 & 6.968945\\
\bottomrule
\end{tabularx}

{\small\hint{Reproduced by final/src/block1\_own\_objective.py: solves build\_b2\_model(d) exactly for the assignment, re-solves SA(seed=0) for the other, and evaluates CP1's own $F(x)=\sum_j n_j(x)u_j(x)^2$ on each resulting assignment directly from the same cpu\_req/cpu\_cap data. Source: results/final\_block1\_own\_objective.json.}}

\subsection*{Conclusion}
The order is \textbf{not} preserved. Under the reference objective the two solutions are exactly tied: both are optimal, at 5.268322. Under our own Checkpoint~1 objective they are not tied at all: the approximate method's assignment scores 6.968945 against the exact solver's assignment's 7.530112 -- an 8.1\% difference that the reference model is structurally blind to, since both assignments look identical to it. So a tie under the model everyone is graded on hides a real, non-trivial preference under the model our team actually started from: if Checkpoint~1's own objective were still the one being optimized, our team should prefer the kind of assignment the approximate methods happen to produce over the one the exact solver on the reference model hands back, even though the reference model itself cannot tell them apart. Both, in turn, cost noticeably more under Checkpoint~1's objective (7.53 and 6.97) than Checkpoint 1's own hand-built witness assignment scored under that same objective (5.075, Appendix C of that report) -- even though that witness was only a greedy construction improved by single-VM moves, not an optimized solution, which shows neither reference-optimal assignment is even close to what Checkpoint 1's own model would call good. The mechanism is the objective change identified in Table 1: the reference model minimizes total latency summed once per server, while Checkpoint 1 minimizes total latency weighted by how many VMs are exposed to it, so a placement that concentrates VMs on a few low-$t_l$ servers can be free under the reference model while being expensive under Checkpoint 1's $n_j(x)u_j(x)^2$ term. The two models are not interchangeable summaries of the same quality, and a result reported under one should never be quoted as if it were optimal, or even good, under the other -- which is also why every number elsewhere in this report (Blocks 2-4) is explicitly reported under the reference model only, never mixed with Checkpoint 1's.

\section{Comparison of the methods of the team}

{\small
\begin{tabular}{@{}p{3.0cm}rrrrccrr@{}}
\toprule
Method & best & mean & worst & std & feas. & time, s & gap$_{\mathrm{ref}}$, \% & gap$_{\mathrm{LP}}$, \%\\
\midrule
GA, penalty (GA-A/GA-P)          & 5.282228 & 5.297992 & 5.312343 & 0.009793 & 10/10 & 12.16  & 0.563 & 1.667\\
SA, transfer+swap (SA-I/SA-T)    & 5.268322 & 5.268322 & 5.268322 & 0.000000 & 10/10 & 0.22   & 0.0000 & 1.113\\
SA, transfer (SA-solo)           & 5.276919 & 5.283488 & 5.295145 & 0.006855 & 10/10 & 0.15   & 0.288 & 1.397\\
Random search                    & 5.549270 & 5.623761 & 5.717093 & 0.051770 & 10/10 & 14.39  & 6.747 & 7.363\\
Exact solver (HiGHS, CP3 model)  & 5.268322 & --       & --       & --       & proven & 1.3    & 0.0000 & 1.113\\
\bottomrule
\end{tabular}}

\[
\mathrm{gap}_{\mathrm{ref}}=100\cdot\frac{z-\zref}{\zref},\qquad
\mathrm{gap}_{\mathrm{LP}}=100\cdot\frac{z-\zLP}{z}.
\]
{\small\hint{Numbers reproduce from results/Zhussupbekova\_runs.csv (GA and the random-search row, the GA pair's baseline), results/tilesbay\_runs.csv (SA, transfer+swap), results/noken\_runs.csv (SA, transfer; its own random baseline, mean 5.585015, leads to the same conclusion) (Checkpoint 4, budget 100\,000, seeds 0--9) and results/final\_exact\_m35.json (this checkpoint's exact-solver rerun on the original instance, CP3 model, HiGHS, proven in 1.3s). $\zref=5.268322$, $\zLP=5.209695$ (task file).}}

\subsection*{Figure F1}
\includegraphics[width=0.95\textwidth]{F1_final.png}

\subsection*{Analysis}

\textbf{Which method is better.} The SA pair (transfer+swap) is best: all ten seeds reach $\zref=5.268322$, the value the exact solver proves optimal (gap$_{\mathrm{ref}}$ 0.000\%, std 0). SA-solo (transfer only) is second (mean 5.283488, gap$_{\mathrm{ref}}$ 0.288\%), GA third (mean 5.297992, gap$_{\mathrm{ref}}$ 0.563\%), random search last (mean 5.623761, gap$_{\mathrm{ref}}$ 6.747\%, more than ten times the gap of GA).

\textbf{Difference against the spread between seeds.} SA pair vs GA: the difference of the means (0.0297) is three times GA's standard deviation (0.0098), and the ranges do not overlap -- GA's best seed (5.282228) is worse than every SA-pair seed, so this difference is real, not seed luck. SA-solo vs GA is closer: the difference of the means (0.0145) is about 1.5 GA standard deviations and 2 SA-solo standard deviations (0.0069), and the ranges overlap (GA best 5.282228 $<$ SA-solo worst 5.295145). SA-solo is better in 87.5\% of the 100 seed pairs: better on average, not on every seed.

\textbf{Figure F1.} GA drops fastest and leads up to about 45\,000 evaluations, then stalls near 5.30. Both SA variants start slowly, because at high temperature they accept many worsening moves, pass GA at about 45\,000--50\,000 evaluations and reach their final values by about 60\,000 evaluations, so the budget of 100\,000 is not the limiting factor for SA.

\textbf{Mechanism: neighbourhood and acceptance, supported by the Checkpoint 4 checks of three members.} SA does not search its neighbourhood exhaustively: each evaluation proposes one random move. What makes it work is the acceptance rule. In Rakhima's (SA-solo) check, SA accepted 52--55\% of the worsening moves in the first 5\% of the budget on every seed (seed 3: 1560/2913) and 0\% in the last 5\% on nine seeds (0.1\% on seed 6; seed 3: 0/4613) (\texttt{results/noken\_analysis.json}): early it escapes local optima, at the end it is a local search that stops in a local optimum of its own neighbourhood.
\begin{itemize}[leftmargin=*,itemsep=1pt]
\item \emph{The neighbourhood explains SA pair vs SA-solo.} At SA-solo's worst seed (seed 3, cost 5.295145) there are 0 improving transfers out of 315 but 7 improving swaps out of 532. The best swap, VM3 (server 9) $\leftrightarrow$ VM32 (server 8), lowers the cost by 0.0221, but it cannot be done as two transfers: moving VM3 first raises the cost by 0.0335 (rejected at $T\approx0.012$), moving VM32 first is infeasible. Tilesbay's (SA-I) check found 0 improving transfers and 0 improving swaps at the final solution (\texttt{results/tilesbay\_analysis.json}): with the swap move added, SA ends in a local optimum of the larger neighbourhood, which here is $\zref$.
\item \emph{No local polishing explains GA.} Danial's (GA-P) worst-seed check (seed 5, cost 5.312343, \texttt{results/kabdulla\_analysis.json}) found 1 improving transfer out of 315 and 7 improving swaps out of 538 next to GA's solution, and the population collapsing from 30 distinct genotypes to 8, with the last improvement at evaluation 41\,114. Crossover of near-identical parents produces almost nothing new and GA has no step that polishes the best solution, so it stops in a point that is not even a local optimum -- the flat GA curve in F1.
\end{itemize}

\textbf{Distance to $\zLP$.} No method can close gap$_{\mathrm{LP}}$: SA pair's gap$_{\mathrm{LP}}$ (1.113\%) equals the integrality gap of the relaxation, $100(\zref-\zLP)/\zref=1.113\%$, because $\zref$ is proven optimal. The remaining distance to $\zLP$ is the weakness of the LP bound, not a search error of any method.
\section{Scaling: where does each method break down?}

All runs below use the same budget and parameters as Checkpoint 4: 100{,}000 evaluations, seeds 0-9, and for SA/SA-solo the same tuned $(T_0,T_{\mathrm{end}})$ pairs as Checkpoint 4 (no re-tuning at the scaled sizes); GA and random search also keep their Checkpoint 4 settings unchanged. The exact solver uses the team's own Checkpoint 3 reference model, \texttt{build\_b2\_model(d)}, built directly on each scaled instance, solved with HiGHS (\texttt{scipy.optimize.milp}) under a 120-second cap.

\subsection*{Table 3}
{\small
\begin{tabular}{@{}lp{3.8cm}rrp{2.0cm}p{2.2cm}r@{}}
\toprule
Size & Method & Value & Best known & Gap to best known, \% & Solver gap, \% (exact only) & Time, s\\
\midrule
$M{=}100$ & exact (HiGHS)            & 9.176682 & 9.176682 & 0.000 & 0.000 & 14.7\\
          & SA, own start            & 9.176682 & 9.176682 & 0.000 & -- & 0.33\\
          & SA, LP-rounded start     & 9.176750 & 9.176682 & 0.001 & -- & 0.31\\
          & SA-solo, own start       & 9.188359 & 9.176682 & 0.127 & -- & 0.10\\
          & SA-solo, LP-rounded start& 9.193204 & 9.176682 & 0.179 & -- & 0.11\\
          & GA, penalty              & 9.229146 & 9.176682 & 0.572 & -- & 92.7\\
          & random search            & 10.817779& 9.176682 & 17.88 & -- & 150.1\\
\addlinespace
$M{=}200$ & exact (HiGHS, not proven)& 24.945887& 24.875537& 0.283 & 0.953 & 120.0\\
          & SA, own start            & 24.884626& 24.875537& 0.037 & -- & 0.38\\
          & SA, LP-rounded start     & 24.884657& 24.875537& 0.037 & -- & 0.40\\
          & SA-solo, own start       & 24.975725& 24.875537& 0.403 & -- & 0.19\\
          & SA-solo, LP-rounded start& 25.018192& 24.875537& 0.573 & -- & 0.17\\
          & GA, penalty              & 25.293901& 24.875537& 1.682 & -- & 107.0\\
          & random search            & 35.851362& 24.875537& 44.12 & -- & 170.8\\
\addlinespace
$M{=}400$ & exact (HiGHS, not proven)& 55.339800& 54.924503& 0.757 & 1.251 & 120.1\\
          & SA, own start            & 54.945982& 54.924503& 0.039 & -- & 0.43\\
          & SA, LP-rounded start     & 54.947637& 54.924503& 0.042 & -- & 0.40\\
          & SA-solo, own start       & 55.193688& 54.924503& 0.490 & -- & 0.36\\
          & SA-solo, LP-rounded start& 55.179680& 54.924503& 0.464 & -- & 0.35\\
          & GA, penalty              & 55.985455& 54.924503& 1.932 & -- & 175.6\\
          & random search            & 105.306838$^*$& 54.924503& 91.73 & -- & 329.0\\
\bottomrule
\end{tabular}}

\noindent{\small $^*$random search was infeasible on 2 of its 10 seeds at $M{=}400$ (8/10 feasible); the value shown is the mean over the 8 feasible runs only, so it understates how often the method fails outright at this size.}

\hint{All ``Value'' entries are the mean over 10 seeds (feasible runs only, see the random-search caveat above); Time is the mean wall-clock seconds per run, recomputed directly from the per-seed `seconds` column of results/final\_scaling\_runs.csv (an earlier draft wrongly rounded all SA/SA-solo times to 0.0 -- fixed here per Tamerlan's review). Source: results/final\_scaling\_runs.csv (scaling matrix) and results/final\_exact\_scaling.json (exact rows).}

\subsection*{On the start argument}
GA and random search are unaffected by the \texttt{start} argument by construction: both always begin from their own fresh random draw (GA's population, random search's first candidate), exactly as Checkpoint 4 required, and the code never reads \texttt{start}. So for these two methods, "own start" and "LP-rounded start" are literally the same run -- a legitimate instance of "nothing," stated here rather than wasting the budget rerunning something that cannot differ. SA and SA-solo do use \texttt{start}: our own rounding rule (documented in \texttt{final/src/build\_rounded\_start.py}) takes the LP relaxation of the team's reference model (\texttt{src/model/build.py}, same model as the exact solver) on the scaled instance and, for each VM, assigns it to the zone-allowed server with the largest LP fraction $y_{k,l}$, breaking ties toward the currently least-loaded allowed server. The rounded start is capacity-feasible at $M{=}100$ but not at $M{=}200$ or $M{=}400$ (a few servers come out slightly over on memory; both SA variants handle an infeasible supplied start by searching from it until a feasible state is found, well inside the budget). The rounded start changes SA's and SA-solo's final cost by at most 0.1\% at any size tested -- a real but small effect, smaller than the gap between SA and SA-solo themselves.

\subsection*{Figure F2}
\includegraphics[width=0.85\textwidth]{F2_final.png}

\subsection*{Analysis}

The exact solver proves optimality at $M{=}100$ in 14.7s and does not prove it at $M{=}200$ or $M{=}400$ within the 120s cap (solver gaps 0.953\% and 1.251\%, both honestly reported from the solver's own best bound, not read off the word "Optimal"). So the break point for this storyline and this model sits between $M{=}100$ and $M{=}200$.

At both sizes where the exact solver fails to prove optimality, SA's own-start result is already \emph{below} the exact solver's own 120-second incumbent: at $M{=}200$, SA reaches 24.884626 against the exact solver's 24.945887; at $M{=}400$, SA reaches 54.945982 against 55.339800. So the approximate method overtakes the time-limited exact solver's incumbent at exactly the sizes where the exact solver stops proving optimality -- not before, since at $M{=}100$ the exact solver already matches SA exactly (both find 9.176682, SA with zero variance).

What three sizes cannot show: whether this crossover is a one-time effect of $M{=}200$ being just past the exact solver's practical limit, or whether the margin between SA and the exact incumbent keeps growing as $M$ grows further -- that needs more than two post-breakdown points to see a trend, and all three sizes used the same 120-second cap, so we cannot tell whether a longer time limit would let the exact solver catch back up at $M{=}200$ or $M{=}400$. An additional experiment that would help: rerun the exact solver at $M{=}200$ with a much longer cap (say 1200s) to see whether its incumbent keeps improving fast enough to overtake SA again, or whether it has effectively stalled on its bound.

GA and random search scale the same way they did at the team's own size: GA stays feasible on every seed at all three sizes but its gap to best known grows with $M$ (0.57\% at 100, 1.68\% at 200, 1.93\% at 400), and it costs far more wall-clock time than SA or SA-solo for it (93-176s mean per run here, against well under one second for either SA variant) without ever beating them on cost. Random search degrades faster still: its gap to best known grows from 17.9\% at $M{=}100$ to 44.1\% at $M{=}200$ to 91.7\% at $M{=}400$, and at $M{=}400$ it is no longer reliably feasible at all, failing on 2 of its 10 seeds (the value reported there is the mean over the 8 feasible runs, so the true gap including failures is worse than the table shows). Both confirm the same scaling story as Table 2 at a larger size: SA's local-move neighbourhood is the only one of the four methods that keeps up with the exact solver as $M$ grows.

\section{Method-selection memo}

\emph{To: the owner of the cloud placement problem. From: Team T12. Re: which optimizer to run.}

\begin{tabularx}{\textwidth}{@{}lXXX@{}}
\toprule
 & Own instance ($M{=}35$) & Medium scaled ($M{=}200$) & Large scaled ($M{=}400$)\\
\midrule
Problem & Place 35 VMs on 10 servers under zone + capacity limits, minimize total latency. & 200 VMs, 57 servers, same rules. & 400 VMs, 114 servers, same rules.\\
Method & Exact solver (HiGHS) & SA, transfer+swap, own start (SA-I/SA-T) & SA, transfer+swap, own start (SA-I/SA-T)\\
Result and its level & 5.268322, proven optimal (gap 0.000\%) & 24.884626, gap to best known 0.037\% & 54.945982, gap to best known 0.039\%\\
Check & Solver's own best bound meets the solution (gap $\le10^{-6}$); cross-checked against SA's independent result, identical to 6 decimals at $M{=}35,100$. & Checker (\texttt{evaluate\_B}) re-verifies feasibility and cost on every reported run; 10/10 seeds feasible, std 0.004 (own start). & Same checker re-verification; 10/10 feasible, std 0.009.\\
Risk & None at this size -- exact and SA agree. & Exact solver cannot prove optimality in 120s at this size, so "0.037\% gap" is relative to the best value found so far, not a mathematical certificate. & Same caveat, larger: exact solver's own gap is 1.25\% here, so SA's advantage over it could narrow or widen with more solver time -- untested beyond 120s.\\
\bottomrule
\end{tabularx}

\subsection*{Recommendation for the real situation}
For the team's own instance size, either method works and the exact solver is actually preferable: it is fast (1.3s) and gives a certified optimum, so there is no reason to settle for a heuristic. Once the problem grows past roughly 100 VMs, switch specifically to \textbf{SA, transfer+swap (the pair's tuned variant, SA-I/SA-T)}, own start -- not SA-solo: at every tested size, the pair's variant beats SA-solo (gap to best known 0.037\% vs.\ 0.403\% at $M{=}200$; 0.039\% vs.\ 0.490\% at $M{=}400$, Table 3). It already matches or beats the exact solver's own 120-second incumbent at $M{=}200$ and $M{=}400$, runs in well under a second, and its result is checker-verified and reproducible to the seed.

GA is not recommended at any size tested here: it is consistently worse than SA on cost, and its runtime grows with size (mean 92.7s at $M{=}100$, 107.0s at $M{=}200$, 175.6s at $M{=}400$; range 53-195s across all individual runs) against SA's well under one second regardless of size -- for no accuracy advantage. Its only edge would be on problems where SA's two-move neighbourhood (transfer, swap) is a poor fit, which this storyline does not show.

Random search is worse again on both axes: cost (17.9\% gap to best known at $M{=}100$, 44.1\% at $M{=}200$, 91.7\% at $M{=}400$ -- growing fastest of any method as $M$ grows) and reliability (10/10 feasible at $M{=}100$ and $M{=}200$, but only 8/10 at $M{=}400$ -- 2 of its 10 seeds fail to find any feasible placement at all within the 100{,}000-evaluation budget at that size). It should not be used past the smallest instance sizes tested here.

\subsection*{What would change the recommendation}
A fill ratio close to 100\% (CPU or memory) would change this: SA's exhaustive local-move search and GA's survive-the-penalty approach both rely on there being room to place VMs without violating capacity almost everywhere, which is why random search is already 10/10 feasible here. If capacity were tight enough that most assignments were infeasible, the zone-respecting-but-capacity-blind operators both methods use (SA's transfer/swap, GA's crossover/mutation) would spend much more of their budget rejecting or penalizing infeasible moves, and a method with an explicit repair step could become competitive again. Similarly, if the exact solver's 120s incumbent kept improving fast in the untested longer-time experiment above, the recommendation to switch away from it at $M{=}200$ would need revisiting.

\appendix
\section{Authorship}
\begin{tabularx}{\textwidth}{@{}XXXXl@{}}
\toprule
Block & Wrote the code & Ran the experiment & Wrote the text & Signature\\
\midrule
1. Model & & & & \\
2. Comparison & & & & \\
3. Scaling & & & & \\
4. Memo & & & & \\
\bottomrule
\end{tabularx}

\section{Declarations}
External code and AI-assistant use are declared in \texttt{README.md}, section ``Final'', one table per author, in the format of Checkpoint 4.

\end{document}
